\documentclass[10pt,a4paper]{amsart}
\usepackage[margin=19mm]{geometry}
\usepackage{amsmath,amssymb,mathtools}
\usepackage[scr=rsfs,cal=euler]{mathalfa}
\usepackage{xfrac}
\usepackage[unicode,hidelinks,bookmarks=false]{hyperref}
\newcommand{\ie}{i.e.\ }
\newcommand{\cf}{cf.\ }
\newcommand{\resp}{respectively\ }
\newcommand{\Subsection}{\subsection}
\providecommand{\nobreakdash}{\nobreak\hbox{-}}
\setlength{\emergencystretch}{2em}
\multlinegap0pt
\usepackage{amssymb}
\usepackage{float}
\usepackage{enumerate}
\usepackage{bm}
\usepackage{xcolor}
\usepackage[normalem]{ulem}
\newtheorem{thm}{Theorem}
\newtheorem{lem}{Lemma}
\usepackage{cleveref}
\crefname{thm}{Theorem}{Theorems}
\crefname{lem}{Lemma}{Lemmas}
\newcommand{\Ann}{{\rm Ann}}
\DeclareMathOperator{\F}{\mathbb{F}}
\DeclareMathOperator{\Z}{\mathbb{Z}}
\newcommand{\soc}{{{\rm soc}}}
\title{Perfect state transfer on \\ gcd-graphs over a finite Frobenius ring}
\author{Tung T. Nguyen, Nguyen Duy T\^an}
\date{Source: arXiv:2504.00404v2 (CC BY 4.0)}
\begin{document}
\maketitle

\noindent\emph{Source and licence.} Perfect state transfer on \\ gcd-graphs over a finite Frobenius ring. Version arXiv:2504.00404v2 (CC BY 4.0), distributed by arXiv under the Creative Commons Attribution 4.0 International licence, http://creativecommons.org/licenses/by/4.0/, as recorded in the arXiv licence metadata for this exact version and shown on the abstract page https://arxiv.org/abs/2504.00404v2. Full licence notice: \texttt{licenses/arxiv-2504.00404-CC-BY-4.0.txt}.\par\smallskip

\noindent\textit{Editorial scope.} Counted unit: the article's thm thm:unitary (source0.tex lines 769--776). The article's complete original proof of it is delivered with it (source0.tex lines 778--813, one \texttt{proof} environment of 36 lines). No mathematics is written, repaired or re-derived: the statement and the proof are the article's own lines, in the article's own notation, and no retained line is substituted or re-flowed.  Retained uncounted as the source-local context the proof needs: lines 324--327; lines 364--366; lines 377--379; lines 544--554.  Nothing was omitted from any retained span, so the package carries no omission record.  No retained line contains a citation, so the wrapper carries no bibliography and no bibliography entry.  No retained line contains a figure environment, an image-inclusion command or drawing code, so no image is delivered and the package declares no asset dependency.  Editorial formatting only, in addition: an AMS \texttt{amsart} preamble that restates the article's own declarations and macros used by the retained lines, namely the environments \texttt{thm}, \texttt{lem} on separate counters, together with the standard packages those lines need.  The retained environments print in this document as Theorem 1, Lemma 1 in order of appearance, whereas the article prints the same items with the numbering of its own full document.  The complete native source of the article as distributed in the arXiv e-print for this exact version is bundled unchanged as \texttt{sources/175750/source0.tex}.
\begin{thm} \label{prop:condition}
    There exists perfect state transfer from $0$ to $s$ at time $t$ if and only if for all $r_1, r_2 \in R$
    \[ (\lambda_{r_1}-\lambda_{r_2}) \frac{t}{2 \pi} + \frac{\psi(s(r_1-r_2))}{n} \equiv 0 \pmod{1}.\]
\end{thm}

\begin{lem} \label{lem:local}
    Let $(S, \mathfrak{m})$ be a local finite Frobenius ring such that the residue field of $S$ is $\F_2.$ Then, there exists a unique $e \in S$ such that $\Ann_{R}(e)=\mathfrak{m}$. Furthermore, if $\psi\colon S \to \Z/n$ is any non-degenerate linear functional then $\dfrac{\psi(e)}{n} \equiv \dfrac{1}{2} \pmod{1}.$ 
\end{lem}

\begin{thm} \label{prop:upper_bound}
Let $R$ be a finite Frobenius ring. Suppose that $R$ has the following Artin-Wedderburn decomposition: $R = (\prod_{i=1}^d S_i) \times R_2$. Here, $(S_i, \mathfrak{m_i})$ represents all local factors of $R$ whose residue fields are $\mathbb{F}_2$. For each $1 \leq i \leq d$, let $e_i$ be the unique minimal element of $S_i$ as defined in \cref{lem:local}. If there exists PST between $0$ and some $s \in R$, then $s$ must be of the form $(a_1, a_2, \ldots, a_d, 0)$, where each $a_i$ is $0$ or $e_i$. In particular, if $R$ is a local ring, then $s=e$, where $e$ is the minimal element of $R$. 
\end{thm}

\begin{lem} \label{lem:secondary_minimal}
\noindent
Let $a \neq e$ be an element of $R$. The following statements are equivalent. 
\begin{enumerate}
\item $Ra$ is a minimal ideal containing $Re.$
    \item  $xa \in Re$ for all $x \in \mathfrak{m}.$
    \item $a \in \Ann_{R}(\mathfrak{m^2}).$
    \item $a \in \soc^2(R).$
    \item $|R/\Ann_{R}(a)| = 4.$ 
\end{enumerate}
\end{lem}

\begin{thm} \label{thm:unitary}
    Suppose that there exists PST on $G_{R}(1)$ with $R=S_1 \times \prod_{i=1}^r \F_{q_i}$. Then the following conditions must hold. 
\begin{enumerate}
    \item $S_1 \in \{\F_2, \Z/4, \F_2[x]/x^2 \}.$
    \item $q_i \equiv 0 \pmod{4}$ for all $i.$
\end{enumerate}
Conversely, suppose that these conditions hold. Then there exists PST between $(0,0)$ and $(e_1, 0)$ at time $\frac{\pi}{2}.$ Here $e_1$ is the minimal element of $S_1.$
\end{thm}

\begin{proof}
    Let us prove the necessary conditions. Suppose that there exists PST between $0$ and some $s \in R.$ By \cref{prop:upper_bound}, $s$ must equal to $(e_1, 0, \ldots, 0).$ In this case, the conditions in \cref{prop:condition} can be described explicitly as 
\begin{equation} \label{eq:conditions}
    \begin{cases}
    (\lambda_{(x_1, a_1, \ldots, a_r)} - \lambda_{(x_2, b_1, \ldots, b_r)}) \frac{t}{2 \pi} + \frac{1}{2} \in \Z  \text{if } (x_1-x_2) \not \in \mathfrak{m}_1 \\ 
    (\lambda_{(x_1, a_1, \ldots, a_r)} - \lambda_{(x_2, b_1, \ldots, b_r)}) \frac{t}{2 \pi}  \in \Z  \text{if  } (x_1-x_2)  \in \mathfrak{m}_1.
\end{cases} 
\end{equation}
This follows from the fact that $\frac{\psi((e_1,0, \ldots, 0))}{n}= \frac{1}{2}$ as shown in \cref{lem:local}. We also note that 
\begin{align*} 
\lambda_{(x_1, a_1, \ldots, a_r)} &= \frac{\varphi(S_1)}{\varphi(S_1/\Ann_{S_1}(x_1))} \mu(S_1/\Ann_{S_1}(x_1)) \prod_{i=1}^r \frac{\varphi(\F_{q_i})}{\varphi(\F_{q_i}/\Ann_{\F_{q_i}}(a_i))} \mu(\F_{q_i}/\Ann_{\F_{q_i}}(a_i)) \\ 
&= \frac{\varphi(S_1)}{\varphi(S_1/\Ann_{S_1}(x_1))} \mu(S_1/\Ann_{S_1}(x_1)) \prod_{a_i =0} (q_i-1) \prod_{a_i \neq 0 }(-1) \\
&= \frac{\varphi(S_1)}{\varphi(S_1/\Ann_{S_1}(x_1))} \mu(S_1/\Ann_{S_1}(x_1)) (-1)^r \prod_{a_i =0} (1-q_i).
\end{align*}



First, we claim that $S_1 \in \{\F_2, \Z/4, \F_2[x]/x^2 \}.$ If $S_1=\F_2$ then we are done. Therefore, to avoid tautology, we will assume that $S_1 \neq \F_2.$ We first claim that $\mathfrak{m}_1=S_1e_1$. In fact, suppose that it is not the case. Then we can find $x \in \mathfrak{m}_1 \setminus Re$ (since we are assuming that $S_1 \neq \F_2$, $S_1e_1 \subset \mathfrak{m}_1).$ We then have $\Ann_{S_1}(x) \neq \mathfrak{m}_1$ and hence $\lambda_{(x,0,\ldots,0)}=0.$ Similarly, we also have $\lambda_{(1,0,\ldots, 0)}=0$. These two equalities contradict the first condition. We conclude that $\mathfrak{m}_1 = S_1 e_1.$ This would imply that $\mathfrak{m}_1^2=0$ and hence $1 \in \Ann_{S_1}(\mathfrak{m}_1^2).$ By \cref{lem:secondary_minimal}, we conclude that $|S_1|=4$ and hence $S_1 \in \{ \Z/4$, $\F_2[x]/x^2 \}.$ 

Next, we claim that $q_i \equiv 0 \pmod{4}$ for each $1 \leq i \leq r.$

\textbf{Case 1.} Let's first consider the case $S_1=\F_2.$ We have $\lambda_{(1, 1, 1,\ldots, 1)}=(-1)^{r+1}$,  $\lambda_{(0,1,\ldots, 1)} =(-1)^r$, and $\lambda_{(1, 0, 1, \ldots, 1)}=(-1)^{r+1}(1-q_1).$ The first condition then implies that $2(\frac{t}{2 \pi})+\frac{1}{2} \in \Z.$ The second condition then implies that $q_1 \frac{t}{2 \pi} \in \Z.$ This condition necessarily implies that $q_1$ must be even. By the same argument, $q_i$ is even for all $1 \leq i \leq r.$

\textbf{Case 2.} Let us consider the case $S_1 \in \{\Z/4, \F_2[x]/x^2 \}.$ In this case $\lambda_{(e, 1, 1,\ldots, 1)}=2(-1)^{r+1}$, $\lambda_{(1,1,\ldots, 1)}=0$,  and $\lambda_{(e,0,1\ldots,1)}= 2(-1)^r (1-q_1).$ The first condition would imply that $2(\frac{t}{2 \pi})+\frac{1}{2} \in \Z.$ The second condition implies that $2(-1)^{r+1}(1-q_1) \frac{t}{2 \pi} \in \Z. $ We conclude that $q_1$ must be even. By the same argument, $q_i$ are even for all $1 \leq i \leq r.$



Conversely, suppose that all the above conditions are satisfied. We claim that there exists PST between $0$ and $(e_1, 0, \ldots, 0)$ at time $\frac{\pi}{2}.$ We consider two separate cases as above. 

In the first case, namely when $S_1 = \F_2$, we can see that $\lambda_{(1,a_1, a_2, \ldots, a_r)} \equiv (-1)^{r+1} \pmod{4}$ and $\lambda_{(0, a_1, \ldots, a_r)} \equiv (-1)^r \pmod{4}$ for every $(a_1, a_2, \ldots, a_r) \in \prod_{i=1}^r \F_{q_i}.$ We can check that the conditions described in \cref{eq:conditions} are satisfied when $t = \frac{\pi}{2}.$


In the second case where $S_1 \in \{\Z/4, \F_2[x]/x^2\}$ we can see that 
\[ \lambda_{(0,a_1, \ldots, a_r)} \equiv \lambda_{(e_1, a_1, \ldots, a_r)} \equiv 2(-1)^r \pmod{4} .\] 
On the other hand, if $u \in S_1^{\times}$ then $\lambda_{(u_1, a_1, \ldots, a_r)} \equiv 0 \pmod{4}$. These congruences show that all conditions described by \cref{eq:conditions} are met when $t = \frac{\pi}{2}.$
\end{proof}

\end{document}
